\chapter*{Abstract}
\addcontentsline{toc}{chapter}{Abstract}
\label{chap:abstract}

Email marketing remains one of the most cost-effective digital communication channels for businesses, yet campaign design continues to rely heavily on A/B testing --- a static, traffic-intensive method that is slow to adapt and incapable of personalizing decisions to individual recipients in real time. This thesis investigates whether AI-based adaptive optimization methods can systematically outperform traditional A/B testing in the context of email campaign design.

We propose and evaluate a unified comparative framework covering four approaches: (1) classical A/B testing as a baseline, (2) Bayesian inference with Thompson Sampling, (3) contextual multi-armed bandits conditioned on user segment features, and (4) reinforcement learning applied to multi-step drip campaign sequences. The comparison is conducted using anonymized historical campaign data from ProvenExpert, a German B2B review platform, and evaluated across click-through rate, open rate, conversion rate, and convergence speed within user segments defined by device type, geographic region, and past engagement behavior.

The experimental results demonstrate that Bayesian bandit methods converge to optimal variant selection significantly faster than A/B testing while maintaining statistical reliability. Contextual bandits achieve measurable personalization gains across user segments, and reinforcement learning shows promise for multi-step campaign sequences, though it requires substantially more data to converge stably.

This thesis contributes a structured comparative framework for adaptive email campaign optimization and provides practical guidelines for marketing teams seeking to move beyond static experimentation toward continuously learning, data-driven campaign systems.

\vspace{1.5cm}
\noindent\textbf{Keywords:} email marketing, A/B testing, Bayesian inference, Thompson Sampling, multi-armed bandits, contextual bandits, reinforcement learning, campaign optimization, click-through rate
